% sections/benchmark_design.tex --- Benchmark Design
\section{Benchmark Design}
\label{sec:benchmark}

This section fixes the \bench{} protocol: the tasks that submissions
target, the splits over which they are evaluated, the metrics
reported, and the reproducibility envelope under which a
submission's numbers are accepted.  All elements of the protocol are
testable from the released artefacts alone, without requiring
any private data or any non-public tooling.

\subsection{Tasks}
\label{sec:benchmark:tasks}

\bench{} defines three edge-level prediction tasks over
\texttt{graphs/je\_network.csv}:

\begin{itemize}
  \item \textbf{T1: Binary anomaly classification.}  Predict
    \texttt{is\_anomaly} for each edge in the test split.  The
    positive class includes all five categories
    (fraud, error, process, statistical, relational).  T1 is the
    headline task and the one against which leaderboard rank is
    computed.
  \item \textbf{T2: Multi-class anomaly classification.}  Predict
    one of $\{$\texttt{none}, \texttt{fraud}, \texttt{error},
    \texttt{process}, \texttt{statistical}, \texttt{relational}$\}$
    for each edge.  T2 supports per-category analysis and is the
    natural target for class-conditional fraud-detection methods.
  \item \textbf{T3: Rare-typology detection.}  A subset of T2 in
    which only the rarest typology classes (cumulatively the
    bottom $5\%$ of labelled edges) appear in the test split.
    T3 is the hardest task: even strong baselines on T1 and T2
    fall back close to base-rate performance under T3.
\end{itemize}

For each task the prediction is per-edge (i.e.\ per row of the
edge-list CSV); models that operate on the entry table, on the
document chain, or on alternate graph projections are welcome,
provided the final output is an edge-level score.

\subsection{Splits}
\label{sec:benchmark:splits}

Three reproducible splits are released with \bench{}.  They are not
mutually exclusive: a method's performance on each is reported
separately, and the leaderboard records all three.

\begin{itemize}
  \item \textbf{Temporal.}  The last three fiscal months of the
    generated period are held out as test; the preceding nine
    months form the training pool, with the last $20\%$ of the
    training pool reserved as a validation slice.  This split
    measures generalisation across time and is the most common
    real-world deployment scenario.
  \item \textbf{Entity.}  Ten percent of subsidiaries
    (sampled uniformly without replacement) are held out as test;
    all postings of held-out subsidiaries are excluded from
    training.  This split measures generalisation to entities
    unseen during training.
  \item \textbf{Fraud-typology.}  Two of the five anomaly
    categories (sampled at random and fixed in the released
    splits file) are held out as test \emph{positives}; the
    training set sees only the other three categories together
    with negatives.  This split measures generalisation across
    fraud typologies and is the basis for T3.
\end{itemize}

The splits are materialised as CSV files
(\texttt{splits/temporal\_train.csv} etc.\@) listing the
\texttt{edge\_id} of every edge in each fold.  No edge appears in
both the train and test fold of the same split.

\subsection{Metrics}
\label{sec:benchmark:metrics}

All tasks are evaluated under heavy class imbalance (positive rate
on the order of $1$--$5\%$ depending on configuration), which
makes precision-recall metrics more informative than ROC-based ones.
Primary and secondary metrics are reported for every submission:

\begin{itemize}
  \item \textbf{Primary.}  $\AUCPR$ (area under the
    precision--recall curve).  Reported with a $95\%$ bootstrap
    confidence interval.
  \item \textbf{Secondary.}  $\AUCROC$, $F_1$ at the
    threshold that maximises validation $F_1$, precision at
    $k = 50, 100, 500$, and recall at the top $1\%, 5\%, 10\%$ of
    scored edges.
  \item \textbf{Per-typology breakdown.}  For T2 and T3, $\AUCPR$
    is also reported per anomaly category.
\end{itemize}

The leaderboard ranks methods by primary $\AUCPR$ on the temporal
split for T1; the entity- and typology-split numbers are reported
alongside but do not enter the headline rank.

\subsection{Leakage audit}
\label{sec:benchmark:leakage}

A standing risk on any benchmark constructed from a single
generated dataset is that information specific to a held-out
posting leaks into the training pool through a different feature
or row.  \bench{} addresses this risk with an explicit leakage
audit, the script for which is part of the release.  The audit
checks the following:

\begin{itemize}
  \item \emph{Document overlap.}  No edge in the test split shares
    a \texttt{document\_id} with any edge in the train split.
  \item \emph{Predecessor closure.}  No edge in the test split has
    a \texttt{predecessor\_edge\_id} pointing to an edge in a
    different split.  This rule is what enforces document-chain
    integrity at the split boundary.
  \item \emph{Entity-time coverage.}  Under the entity split, no
    held-out subsidiary appears in any aggregated feature
    (consolidation entries, intercompany pairs) that would leak
    test-time data into training.
\end{itemize}

A submission that opts to engineer features from companion files
(e.g.\ document-flow tables, sub-ledger reports) is required to
re-run the audit on its feature-construction script.

\subsection{Reproducibility envelope}
\label{sec:benchmark:repro}

Every submission must record the following metadata in the
leaderboard entry:

\begin{itemize}
  \item the \system{} version (git tag or release number) used to
    produce the data;
  \item the configuration hash (SHA-256 of the canonical YAML);
  \item the ChaCha8 seed under which the dataset was generated;
  \item the model code repository and commit hash;
  \item the random seed(s) used in training.
\end{itemize}

Two submissions whose first three values agree are required to
report identical edge-list contents; this is enforced by a
checksum the audit script computes.  Two submissions whose model
code is identical and whose training seeds match are required to
report identical predictions.

\subsection{Submission protocol}
\label{sec:benchmark:submission}

A submission consists of the following artefacts, packaged in a
single archive:

\begin{itemize}
  \item per-task prediction files
    (\texttt{predictions/T1.csv} etc.), each with two columns
    \texttt{edge\_id} and \texttt{score} (and an additional
    \texttt{predicted\_class} column for T2/T3);
  \item the metadata block listed in
    Section~\ref{sec:benchmark:repro};
  \item a description of the training pipeline of length not
    exceeding two pages, in PDF, sufficient for an independent
    reader to reproduce the result.
\end{itemize}

The release ships a scoring harness
(\texttt{tools/bench\_score.py}) that consumes a submission archive
and emits the full metric table.  Leaderboard entries are accepted
on the basis of harness output; private re-evaluations are not
required, but are encouraged for reviewers who wish to spot-check.
